\section{Key Component Analysis}
\label{sec:components}

This section analyzes six key components in order from core to periphery:
foundation model (corresponding to the CPU) $\to$ KV cache (corresponding to the cache hierarchy) $\to$ context management (corresponding to virtual memory) $\to$ agent runtime (corresponding to the operating system) $\to$ tool bus (corresponding to I/O) $\to$ multi-agent collaboration (corresponding to distributed systems).
For each component, we follow a uniform four-part structure: (1)~the classical computer architecture counterpart, (2)~the corresponding construct in LLM-based systems, (3)~why the two are analogous at the mechanism level, and (4)~the essential differences that the analogy should not obscure.

%% ---------- Foundation Model: Computing Core ----------
\subsection{Foundation Model: Probabilistic Computing Core}
\label{sec:model}

\subsubsection{The Counterpart in LLM-Based Systems}

\textit{Classical counterpart---the CPU:} an execution core that maps input state to output state by running a formally specified instruction set (the ISA) with full \textbf{determinism}---\texttt{ADD R1, R2, R3} executed a million times yields the identical result~\cite{riscv_spec_2026}.

In a model-native computing system, the foundation model (e.g., GPT-4, Claude, LLaMA) plays the role of a \textit{general-purpose execution core}.
Given an input context comprising a system prompt, user messages, tool results, and control constraints, the model produces the next action or output through a forward pass~\cite{vaswani2017attention,brown2020gpt3,llama2023}.

Prompts, tool descriptions, skill files, and project memory files (such as \texttt{AGENTS.md} or \texttt{CLAUDE.md}) can be understood as a form of \textbf{soft instruction interface}: they do not alter the model's physical weights (just as software does not alter the CPU's circuitry), yet they significantly reshape the model's execution path~\cite{openai_codex_agents_md_2026,openai_codex_skills_2026,anthropic_claudecode_memory_2026}.
For example, a carefully authored \texttt{CLAUDE.md} file can cause the same model to exhibit ``Python expert'' behavior in one project and ``Rust expert'' behavior in another, much as loading different programs onto the same CPU yields different behaviors.

SGLang has already moved in this direction: it reformulates complex language-model programs as a collaboration between a structured front-end language and a high-performance runtime, analogous to the transition from assembly language to high-level languages~\cite{sglang2024,sglang_docs_2026}.

\subsubsection{Why the Analogy Holds}

The analogy holds because the two cores share the same relationship between substrate and behavior.
In both, the substrate is fixed while behavior is customized through an external layer---programs for the CPU, prompts and tool schemas for the model---so the same silicon runs a word processor or a simulation just as the same weights write code, analyze data, translate languages, or plan a task.
Both also keep a stable interface while the implementation evolves beneath it.
A CPU exposes its operations through the ISA (arithmetic, branching, memory access); a model exposes its operations through prompt templates and tool schemas (reasoning, generation, tool invocation), and both adhere to the principle of a stable interface over an evolvable implementation.
The x86 architecture has persisted for over four decades while its microarchitecture traveled from single-scalar to superscalar out-of-order execution, much as the GPT-4 API has remained stable across successive weight updates.
Below that interface lies a rich optimization space, and in both cases engineers chase performance without breaking the contract above: CPUs benefit from pipelining, branch prediction, and superscalar issue, while models benefit from FlashAttention (I/O-aware attention computation~\cite{dao2022flashattention}), speculative decoding (using a small model to predict a larger model's output~\cite{leviathan2023speculative}), and continuous batching.

\subsubsection{Essential Differences}

The crux of the difference is \textbf{determinism versus probability}.
A CPU executes a formal instruction set in which every instruction has precisely defined, long-term-stable, fully reproducible semantics.
A foundation model processes a mixture of natural language, code, and structured schemas, producing outputs drawn from a probability distribution.
As Mi et al.\ have argued, drawing on the evolution of computer architecture to inform agent system design requires candid acknowledgment of this fundamental gap between probabilistic and deterministic execution~\cite{mi2025llmagentscomputersystems}.

Consequently, the future ``intelligent ISA'' should not be understood as hard-coding prompts as fixed instructions. Rather, the research agenda should prioritize developing more stable interface layers: tool schemas, constrained output formats, reusable skill packages, task-specific DSLs, and explicit context declarations.
From an architectural perspective, this means the community should evaluate models not only by parameter count but also by \textbf{model programmability}, namely the degree to which upper layers can reliably and predictably control model behavior.

%% ---------- KV Cache: Cache Hierarchy ----------
\subsection{KV Cache: Cache Hierarchy for Intelligent Systems}
\label{sec:kv}

\subsubsection{The Counterpart in LLM-Based Systems}

\textit{Classical counterpart---the cache hierarchy:} L1/L2/L3 caches bridge the processor--memory speed gap by exploiting \textbf{temporal locality}, managed through placement, replacement (e.g., LRU), and write policies~\cite{drepper2007memory}.

During autoregressive inference, generating each new token requires computing attention between that token's Query vector and the Key/Value vectors of \textit{all} preceding tokens.
Recomputing the full prefix KV at every step would incur $O(n^2)$ cost in sequence length.
The KV cache avoids this by storing previously computed Key and Value vectors, reducing each attention step to an $O(1)$ lookup and bringing overall inference cost from $O(n^2)$ to $O(n)$.
This rationale mirrors that of classical caching: \textbf{trade capacity (GPU memory) for latency and bandwidth advantages}~\cite{drepper2007memory,kwon2023pagedattention}.

\paragraph{A KV cache hierarchy is emerging:}
\begin{itemize}[nosep]
\item \textbf{L1---Session-level KV cache:} KV vectors within a single session; lowest latency, capacity bounded by GPU memory.
\item \textbf{L2---Prefix/shared cache:} KV blocks shared across requests for common system prompts or codebase indices, analogous to read-only code pages shared among processes in a CPU~\cite{vllm_prefix_2026,sglang_docs_2026,gim2024promptcache}.
\item \textbf{L3---Semantic cache:} Reusable tool-call results for similar queries, or reusable indices for identical repository structures, producing ``semantic-level'' cache hits.
\end{itemize}

When vLLM introduced PagedAttention, which organizes KV vectors into fixed-size blocks mapped through a block table to non-contiguous physical memory, the analogy ceased to be merely heuristic and became a direct borrowing of virtual-memory page management~\cite{kwon2023pagedattention}.

\subsubsection{Why the Analogy Holds}

The similarity between KV caches and CPU caches operates at the level of engineering mechanisms, and four of those mechanisms line up almost one-to-one.

\paragraph{Shared reliance on temporal locality.}
CPU caches retain recently accessed cache lines because data used recently will be used again soon, and KV caches retain recently computed KV vectors because generating the next token requires attending to all prior tokens.
Both keep frequently reused data in fast storage to avoid expensive recomputation or memory accesses.

\paragraph{Capacity pressure and replacement decisions.}
L1 caches have limited capacity and employ LRU-style replacement policies, and KV caches face analogous GPU memory pressure.
H2O identifies ``heavy-hitter'' tokens (those receiving the highest attention weights) and implements intelligent eviction~\cite{h2o2023}, analogous to identifying ``hot data'' in CPU caches.
StreamingLLM discovered that initial tokens receive disproportionately high attention (the ``attention sink'' phenomenon) and achieved unlimited-length streaming inference by pinning these sink tokens~\cite{xiao2024streamingllm}, which is directly analogous to cache \textit{pinning} in CPU caches.

\paragraph{Bandwidth as the binding bottleneck.}
Gholami et al.\ observe that LLM inference during the decode phase is fundamentally memory-bandwidth-bound, not compute-bound~\cite{gholami2024memorywall}, and Yuan et al.\ systematically characterized this behavior using the Roofline model~\cite{yuan2024llminference}.
Li et al.\ organized KV cache management strategies into three tiers---token-level (which KV entries participate in attention), model-level (architectural optimizations such as Grouped-Query Attention (GQA) and Multi-Query Attention (MQA)), and system-level (memory management and scheduling)~\cite{likvcache2024survey}---mirroring the multi-level approach taken in CPU cache optimization.

\paragraph{Compression as a shared optimization direction.}
KVQuant achieved $8\times$ compression~\cite{hooper2024kvquant}, and KIVI applied asymmetric 2-bit quantization with nearly lossless accuracy and dramatically reduced memory footprint~\cite{liu2024kivi}, paralleling data compression in CPU caches (e.g., IBM MXT) to increase effective capacity.

\subsubsection{Essential Differences}

\paragraph{Boolean versus continuous hits.}
CPU cache hits are \textbf{Boolean}: the address matches or it does not.
In the KV cache, prefix caching does involve exact matching (identical prefix token sequences are reused), but upper-level semantic cache hits are \textbf{continuous-valued}: a similarity score exceeding a threshold is treated as a hit, and different thresholds yield different precision--recall trade-offs.
Classical cache hit-rate models therefore do not transfer directly to semantic caching.

\paragraph{Hardware coherence versus semantic staleness.}
CPU cache coherence is ensured by hardware protocols such as MESI: when data is modified, the corresponding cache lines are marked invalid.
KV and semantic cache ``invalidation'' is far more subtle: tool version changes, code repository commits, and stale indices can all produce hits that return outdated answers, a phenomenon we might call \textit{semantic staleness}.
Recent work has begun exploring learned eviction policies~\cite{lpc2025} and dependency-aware invalidation mechanisms~\cite{multilevelcache2025} to address this challenge.

\paragraph{A steeper capacity--latency gradient.}
In the CPU cache hierarchy, latency increases gradually from L1 to L2 to L3 to main memory on a nanosecond scale.
In the KV cache hierarchy, the latency jump from GPU memory to CPU memory to disk or network spans milliseconds, with no smooth intermediate gradient.
Future cache management will need to co-design traditional cache policies with provenance tracking, TTLs, and hash-based versioning.

Figure~\ref{fig_kv_hierarchy} illustrates the structural correspondence between the classical CPU cache hierarchy and the emerging LLM KV cache hierarchy.
\begin{figure}[htbp]
\centering
\resizebox{\textwidth}{!}{%
\begin{tikzpicture}[
    every node/.style={font=\small},
]

%% Tier y-coordinates (fixed):
%% Tier 1 (fastest): y = 3.5 -- 4.7
%% Tier 2:           y = 2.3 -- 3.5
%% Tier 3:           y = 1.1 -- 2.3
%% Tier 4 (largest): y = 0.0 -- 1.1
%% Half-widths grow downward: 2.1, 2.6, 3.1, 3.6
%% Left column centre: x=0;  Right column centre: x=12

%% ── LEFT: Classical CPU Cache ────────────────────────────────────────────────
% Tier 1: L1
\fill[classical!60, rounded corners=3pt]
    (-2.1, 3.5) rectangle (2.1, 4.7);
\node[text=white, font=\small\bfseries] at (0, 4.22) {L1 Cache};
\node[font=\scriptsize, text=classical!15!white] at (0, 3.88) {$\sim$1\,ns \enspace 64\,KB};

% Tier 2: L2
\fill[classical!42, rounded corners=3pt]
    (-2.6, 2.3) rectangle (2.6, 3.5);
\node[font=\small\bfseries] at (0, 3.02) {L2 Cache};
\node[font=\scriptsize, text=black!60] at (0, 2.68) {$\sim$4\,ns \enspace 512\,KB};

% Tier 3: L3
\fill[classical!25, rounded corners=3pt]
    (-3.1, 1.1) rectangle (3.1, 2.3);
\node[font=\small\bfseries] at (0, 1.82) {L3 Cache};
\node[font=\scriptsize, text=black!55] at (0, 1.48) {$\sim$12\,ns \enspace 8--32\,MB};

% Tier 4: Main Memory
\fill[classical!12, rounded corners=3pt]
    (-3.6, 0.0) rectangle (3.6, 1.1);
\node[font=\small\bfseries] at (0, 0.72) {Main Memory (DRAM)};
\node[font=\scriptsize, text=black!50] at (0, 0.28) {$\sim$60\,ns \enspace 8--64\,GB};

% Left column title
\node[font=\small\bfseries, text=classical!80!black] at (0, 5.05)
    {Classical CPU Cache};

%% ── RIGHT: LLM KV Cache Hierarchy ───────────────────────────────────────────
% Tier 1: Session KV
\fill[modelnative!60, rounded corners=3pt]
    (9.9, 3.5) rectangle (14.1, 4.7);
\node[text=white, font=\small\bfseries] at (12, 4.22) {Session KV Cache};
\node[font=\scriptsize, text=modelnative!15!white] at (12, 3.88) {$<$1\,ms \enspace GPU VRAM};

% Tier 2: Prefix / Shared Cache
\fill[modelnative!42, rounded corners=3pt]
    (9.4, 2.3) rectangle (14.6, 3.5);
\node[font=\small\bfseries] at (12, 3.02) {Prefix / Shared Cache};
\node[font=\scriptsize, text=black!60] at (12, 2.68) {$\sim$1--5\,ms \enspace CPU Memory};

% Tier 3: Semantic Cache
\fill[modelnative!25, rounded corners=3pt]
    (8.9, 1.1) rectangle (15.1, 2.3);
\node[font=\small\bfseries] at (12, 1.82) {Semantic Cache};
\node[font=\scriptsize, text=black!55] at (12, 1.48) {$\sim$10\,ms \enspace SSD / Object Store};

% Tier 4: External Memory
\fill[modelnative!12, rounded corners=3pt]
    (8.4, 0.0) rectangle (15.6, 1.1);
\node[font=\small\bfseries] at (12, 0.72) {External Memory (Vector DB)};
\node[font=\scriptsize, text=black!50] at (12, 0.28) {$\sim$50--200\,ms \enspace Unlimited};

% Right column title
\node[font=\small\bfseries, text=modelnative!80!black] at (12, 5.05)
    {LLM KV Cache Hierarchy};

%% ── Correspondence arrows in the gap (x: 2.1 to 9.9) ────────────────────────
\draw[dashed, gray!55, thick, -{Latex[length=2mm]}]
    (2.2, 4.1) -- (9.8, 4.1);
\node[font=\scriptsize, text=gray!65] at (6, 4.35) {\textit{locality}};

\draw[dashed, gray!55, thick, -{Latex[length=2mm]}]
    (2.7, 2.9) -- (9.3, 2.9);
\node[font=\scriptsize, text=gray!65] at (6, 3.15) {\textit{sharing}};

\draw[dashed, gray!55, thick, -{Latex[length=2mm]}]
    (3.2, 1.7) -- (8.8, 1.7);
\node[font=\scriptsize, text=gray!65] at (6, 1.95) {\textit{similarity}};

\draw[dashed, gray!55, thick, -{Latex[length=2mm]}]
    (3.7, 0.55) -- (8.3, 0.55);
\node[font=\scriptsize, text=gray!65] at (6, 0.80) {\textit{lookup}};

%% ── Capacity / Latency gradient arrow on far left ────────────────────────────
\draw[{Latex[length=2.5mm]}-{Latex[length=2.5mm]}, gray!45, line width=0.6pt]
    (-4.4, 0.1) -- (-4.4, 4.6);
\node[font=\scriptsize, text=gray!55, rotate=90] at (-4.85, 2.35)
    {Capacity $\uparrow$ \quad Latency $\uparrow$};

\end{tikzpicture}
}%
\caption{Structural correspondence between the classical CPU cache hierarchy and the emerging LLM KV cache hierarchy. Both exploit locality and reuse at multiple granularities, with capacity increasing and latency worsening at each lower tier.}
\label{fig_kv_hierarchy}
\end{figure}


%% ---------- Context Management: Virtual Memory ----------
\subsection{Context Management: Virtual Memory with Semantics}
\label{sec:memory}

\subsubsection{The Counterpart in LLM-Based Systems}

\textit{Classical counterpart---virtual memory:} paging gives each process an apparently unbounded address space backed by limited physical memory through \textbf{demand paging}, supported by page tables, TLBs, and replacement policies~\cite{ostep_2023}.

The context window is to a foundation model what main memory is to a process: a high-bandwidth, low-latency, but capacity-limited \textbf{working set}.
Anthropic defines the context window directly as the working memory available to the model during generation.

Context window capacity is expanding rapidly: Gemini~1.5 Pro supports up to 10 million multimodal tokens~\cite{gemini2024}; Ring Attention distributes sequences across multiple devices to achieve nearly unlimited context length~\cite{liu2024ringattention}; positional encoding methods such as RoPE, YaRN, and LongRoPE aim to extend the addressable history range~\cite{su2021roformer,yarn2024,longrope2024}.

Yet a growing body of evidence shows that ``declared support length'' and ``effective working-set capacity'' are not the same thing.
RULER, LongBench~v2, and LOFT provide more realistic stress tests for long-context capability~\cite{ruler2024,longbenchv2_2024,loft2024}.
A machine may advertise 256\,TB of virtual address space while having only 16\,GB of physical memory; true performance depends on whether the working set fits in physical memory at all times.

The more principled design, therefore, should not aim to ``stuff everything into context at once,'' but rather adopt \textbf{tiered management} in the spirit of virtual memory.

\paragraph{An emerging tiered scheme:}
\begin{itemize}[nosep]
\item \textbf{Hot memory (within the context window):} Complete information for the current task, analogous to active pages in physical memory.
\item \textbf{Warm memory (summaries and indices):} Compressed historical context, analogous to pages swapped out but retaining summary metadata.
\item \textbf{Cold memory (external knowledge bases):} Persistent vector stores and document repositories, analogous to data on disk.
\end{itemize}

Figure~\ref{fig:context_tiers} illustrates this three-tier hierarchy alongside its classical virtual-memory counterpart.
% ---------------------------------------------------------------------------
% Context Management Three-Tier Memory Hierarchy
% fig:context_tiers
% ---------------------------------------------------------------------------
\begin{figure}[htbp]
\centering
\providecolor{hotColor}{HTML}{C44E52}
\providecolor{warmColor}{HTML}{E07B39}
\providecolor{coldColor}{HTML}{4A78A8}
\providecolor{classicHot}{HTML}{8B3A8B}
\providecolor{classicWarm}{HTML}{5C7EA8}
\providecolor{classicCold}{HTML}{3A5A8A}
%
\resizebox{0.88\textwidth}{!}{%
\begin{tikzpicture}[
    font=\small, >=Latex,
    tierbox/.style={draw=#1!70!black, fill=#1!18, rounded corners=4pt,
        align=center, font=\small\bfseries, text=#1!85!black, line width=0.7pt},
    classicbox/.style={draw=#1!70!black, fill=#1!12, rounded corners=4pt,
        align=center, font=\small, text=#1!75!black, line width=0.5pt},
    arrowlbl/.style={font=\scriptsize, text=black!55},
]

% ============================================================
%  LEFT SIDE: LLM Context Tiers
% ============================================================
\def\xL{-5.0}

% Hot tier
\node[tierbox=hotColor, minimum width=5.5cm, minimum height=1.5cm]
    (hot) at (\xL, 3.8)
    {Hot Memory\\
     {\normalfont\scriptsize Context window (128K--1M tokens)}\\
     {\normalfont\tiny Latency: $\sim$0\,ms\enspace|\enspace lossless}};

% Warm tier
\node[tierbox=warmColor, minimum width=5.5cm, minimum height=1.5cm]
    (warm) at (\xL, 1.7)
    {Warm Memory\\
     {\normalfont\scriptsize Summaries \& indices (compressed)}\\
     {\normalfont\tiny Latency: $\sim$10--100\,ms\enspace|\enspace lossy}};

% Cold tier
\node[tierbox=coldColor, minimum width=5.5cm, minimum height=1.5cm]
    (cold) at (\xL, -0.4)
    {Cold Memory\\
     {\normalfont\scriptsize External knowledge bases (vector stores)}\\
     {\normalfont\tiny Latency: $\sim$100\,ms--1\,s\enspace|\enspace top-$k$ retrieval}};

% Title
\node[font=\small\bfseries, text=black!70] at (\xL, 5.1) {Model-Native Context Hierarchy};

% Page-in arrows (left side, between tiers)
\draw[-{Latex[length=2mm]}, warmColor!70, thick]
    (warm.north) -- node[arrowlbl, left, xshift=-2pt] {\textit{page-in}} (hot.south);
\draw[-{Latex[length=2mm]}, coldColor!70, thick]
    (cold.north) -- node[arrowlbl, left, xshift=-2pt] {\textit{RAG / page-in}} (warm.south);

% ============================================================
%  RIGHT SIDE: Classical Memory Hierarchy
% ============================================================
\def\xR{5.0}

\node[classicbox=classicHot, minimum width=5.5cm, minimum height=1.5cm]
    (l1) at (\xR, 3.8)
    {Physical Memory (RAM)\\
     {\scriptsize 16--256\,GB; active pages}\\
     {\tiny Latency: $\sim$10--100\,ns\enspace|\enspace lossless}};

\node[classicbox=classicWarm, minimum width=5.5cm, minimum height=1.5cm]
    (l2) at (\xR, 1.7)
    {Swap Space\\
     {\scriptsize Compressed pages / SSD swap}\\
     {\tiny Latency: $\sim$1--10\,ms\enspace|\enspace lossless}};

\node[classicbox=classicCold, minimum width=5.5cm, minimum height=1.5cm]
    (l3) at (\xR, -0.4)
    {Disk / Object Storage\\
     {\scriptsize Files, databases, archives}\\
     {\tiny Latency: $\sim$10--100\,ms\enspace|\enspace lossless}};

\node[font=\small\bfseries, text=black!70] at (\xR, 5.1) {Classical Virtual Memory};

% Page-in arrows (right side)
\draw[-{Latex[length=2mm]}, classicWarm!70, thick]
    (l2.north) -- node[arrowlbl, right, xshift=2pt] {\textit{page-in}} (l1.south);
\draw[-{Latex[length=2mm]}, classicCold!70, thick]
    (l3.north) -- node[arrowlbl, right, xshift=2pt] {\textit{demand paging}} (l2.south);

% ============================================================
%  CENTRE: analogy arrows (short labels only, no overlap)
% ============================================================
\draw[{Latex[length=2mm]}-{Latex[length=2mm]}, thick, black!30, dashed]
    (hot.east)  -- node[font=\tiny, above, text=black!40, yshift=1pt]
    {active pages} (l1.west);
\draw[{Latex[length=2mm]}-{Latex[length=2mm]}, thick, black!30, dashed]
    (warm.east) -- node[font=\tiny, above, text=black!40, yshift=1pt]
    {swap / compressed} (l2.west);
\draw[{Latex[length=2mm]}-{Latex[length=2mm]}, thick, black!30, dashed]
    (cold.east) -- node[font=\tiny, above, text=black!40, yshift=1pt]
    {disk storage} (l3.west);

% ============================================================
%  KEY DIFFERENCE note
% ============================================================
\node[draw=black!25, fill=black!4, rounded corners=4pt,
      font=\scriptsize, align=center, text width=11.5cm,
      inner sep=6pt] at (0, -2.5)
    {\textbf{Key difference:} Virtual memory is \textit{lossless} (exact byte preservation);
     context management is \textit{lossy} (summarization discards detail).
     ``Address translation'' is exact in MMU, but approximate (top-$k$ similarity) in RAG retrieval.};

\end{tikzpicture}}%
\caption{The hot/warm/cold tiered context management hierarchy (left) and its classical virtual-memory analogy (right). Both exploit temporal locality and demand-loading to give consumers the illusion of abundant capacity. The critical distinction is that context management is lossy: summaries and semantic retrieval introduce irreversible information loss absent in classical paging.}
\label{fig:context_tiers}
\end{figure}


MemGPT formalizes this idea as \textit{virtual context management}, explicitly drawing the LLM-as-operating-system analogy and paging information between the context window (``main memory'') and external storage (``disk'')~\cite{memgpt2023}.
LongMem decouples long-term memory from the backbone model entirely~\cite{longmem2023}.
RAG functions as an \textbf{external page-in mechanism}: when the agent encounters a ``page fault,'' it retrieves relevant content from a knowledge base into the context window via semantic search~\cite{lewis2020rag}.
MemoryOS applies OS memory management principles directly to agent memory systems~\cite{memoryos2025}.
MemOS proposes a comprehensive memory operating system architecture~\cite{memos2025}.
HiAgent, inspired by human problem-solving processes, uses sub-goals as memory blocks to hierarchically manage agent working memory~\cite{hiagent2024}.
A-MEM enables agents to organize memory structures autonomously, without predefined schemas~\cite{amem2025}.

\subsubsection{Why the Analogy Holds}

The analogy between context management and virtual memory holds at the \textbf{mechanism level}, and the mechanisms rhyme across three dimensions.

\paragraph{The illusion of abundant space.}
Virtual memory convinces each process that it has a contiguous, large address space, when physical memory is far smaller; virtual context management convinces each agent that it possesses unlimited memory, when the actual context window spans only 128K--1M tokens.
Both realize this illusion through \textbf{demand loading}: only information needed right now occupies ``physical memory'' (the context window), while the rest resides on ``disk'' (external knowledge stores).

\paragraph{Address translation, exact and approximate.}
Virtual memory uses page tables to map virtual addresses to physical addresses; virtual context uses semantic retrieval to map information needs to specific knowledge entries.
The former is an exact match (page number to physical page number); the latter is an approximate match (query intent to top-$k$ relevant documents).

\paragraph{Paging under pressure.}
When physical memory runs low, the OS must decide which pages to swap out (LRU, Clock, and related policies); when the context window fills up, the agent must decide which information to evict (summarization, selective forgetting).
This is precisely the function of the ``context compiler'': the system must determine which state is loaded verbatim, which as a summary, which retains only an index entry, and which is discarded entirely~\cite{contextcompiler2025}.

\subsubsection{Essential Differences}

\paragraph{The key distinction: virtual memory is lossless; context management is lossy.}
OS paging is fully transparent to the application: pages swapped in and out retain every byte unchanged.
Agent context management, by contrast, involves \textbf{semantic compression}: summarization inevitably loses detail, and selection inevitably omits information.
``Address translation'' is also imprecise: vector similarity search returns the ``most relevant'' results, not the ``exact match.''

Semantic virtual memory can therefore never be a lossless paged system.
It is better understood as ``paged memory augmented with semantic summarization and evidence backfill'': information may be distorted after compression, but provenance chains enable verification against original data.
Research from JetBrains suggests that effective context management depends not on ``cramming in more tokens'' but on ``reducing noise and preserving signal''~\cite{jetbrains2025context}, further validating the necessity of lossy compression in this regime.

%% ---------- Agent Runtime: Operating System ----------
\subsection{Agent Runtime: The Intelligent Operating System}
\label{sec:agent}

\subsubsection{The Counterpart in LLM-Based Systems}

\textit{Classical counterpart---the operating system:} system software providing \textbf{virtualization}, \textbf{concurrency}, and \textbf{persistence}---scheduling, isolation, memory management, permissions, I/O, and logging~\cite{ostep_2023}.

If the foundation model is the CPU, then the agent framework increasingly resembles the OS.
An agent runtime must:
\begin{itemize}[nosep]
\item \textbf{Parse objectives:} Decompose user requests into executable subtasks (analogous to program loading and parsing).
\item \textbf{Decide whether to plan:} Route simple requests directly and decompose complex ones into multi-step plans (analogous to fast-path versus slow-path scheduling).
\item \textbf{Load context:} Populate the context window with system prompts, project memory, and conversation history (analogous to loading a process memory image).
\item \textbf{Execute tool calls:} Invoke external tools via protocols such as MCP (analogous to system calls).
\item \textbf{Approve dangerous actions:} Perform permission checks on side-effecting operations such as file writes or code execution (analogous to capability-based access control).
\item \textbf{Spawn sub-agents concurrently:} Dispatch subtasks to independent child agents (analogous to \texttt{fork}).
\item \textbf{Merge intermediate results:} Aggregate sub-agent outputs into a final answer (analogous to inter-process communication and result aggregation).
\item \textbf{Persist state:} Write important state to memory files or knowledge bases (analogous to filesystem writes).
\item \textbf{Recover from failure:} Retry or roll back on errors to maintain task consistency (analogous to transaction recovery).
\end{itemize}

This analogy has already been articulated in the literature.
Ge et al.\ proposed the vision of ``LLM as OS, Agents as Apps''~\cite{ge2023llmos}.
Mei et al.\ built AIOS, a full LLM agent operating system kernel architecture~\cite{mei2024aios}.
Mi et al.\ systematically applied the evolutionary experience of computer architecture to agent system design~\cite{mi2025llmagentscomputersystems}.
Karpathy argued in 2023 that ``LLMs are not chatbots; they are the kernel process of a new operating system,''~\cite{karpathy2023lmos} where the filesystem becomes a vector database and the browser becomes internet search.
This vision is becoming an economic reality: an Augment Code enterprise customer completed a project in two weeks that their CTO had estimated would take four to eight months~\cite{anthropic2026agenticreport}, demonstrating that effective L3 context management combined with L5 orchestration can compress developer onboarding and project delivery from weeks to days.

Industrial practice is evolving in the same direction.
Codex provides subagents, skills, and sandbox-based control surfaces~\cite{openai_codex_overview_2026}; Claude Code offers hooks, MCP integration, memory files, and read-only approval flows~\cite{anthropic_claudecode_overview_2026}.

\subsubsection{Why the Analogy Holds}

Among all the analogies in this paper, the correspondence between the agent runtime and the OS is the deepest, because it operates at the level of \textbf{architectural responsibility}: both take ownership of the same three concerns.

\paragraph{Resource management.}
The OS manages CPU time slices, memory pages, disk space, and network bandwidth; the agent runtime manages model inference quota, context window space, tool-call rate limits, and sub-agent count.
Both must answer the same question: how should limited resources be allocated among competing demands?

\paragraph{Isolation and security.}
The OS isolates processes through address spaces so that one process crash does not affect others~\cite{ostep_2023}, and the agent runtime isolates sub-agents through sandboxes so that one sub-agent's erroneous operation does not compromise the main task~\cite{openai_codex_sandbox_2026,anthropic_claudecode_sandbox_2026}.
Both face the same access-control question: who is allowed to do what?

\paragraph{Standard interfaces.}
The OS standardizes application--hardware interaction through system calls (the POSIX interface); the agent runtime standardizes model--tool interaction through protocols such as MCP~\cite{mcp_intro_2026}.

\subsubsection{Essential Differences}

\paragraph{Unpredictable ``process'' behavior.}
A traditional OS schedules processes whose behavior is defined by deterministic code, so it can accurately predict execution time and resource requirements.
An agent's behavior is determined by probabilistic reasoning: the same task may follow entirely different execution paths across runs, succeeding once and falling into an infinite loop the next, which poses a fundamental challenge for scheduling and resource management.

\paragraph{Fragile shared state.}
When multiple processes share memory, consistency is ensured by hardware cache coherence protocols (e.g., MESI) and software locks, which are precise and verifiable.
When multiple agents share semantic state (e.g., ``an understanding of the same document''), consistency can only be maintained through natural-language communication or external synchronization mechanisms, which are inherently ambiguous and error-prone.

\paragraph{Imprecise interface specifications.}
Every system call in the POSIX interface has precisely defined semantics, parameter types, error codes, and boundary conditions.
An agent's tool-call interface can describe parameters via JSON Schema, but the ``semantic behavior'' of a tool (what it returns under which circumstances) is far from the level of precision expected of a system call.

\textbf{The greatest value of this analogy} lies in reframing many agent problems as OS problems:
\begin{itemize}[nosep]
\item Does multi-agent collaboration resemble multithreading or multiprocessing?~\cite{wu2023autogen}
\item Should tool invocations be treated as system calls?
\item Is an MCP server analogous to a device driver or bus protocol?~\cite{mcp_intro_2026}
\item Why do dangerous commands require capability-based permissions?~\cite{debenedetti2025camel}
\item How can long-running tasks recover via logs and interrupts?~\cite{openai_codex_approvals_2026}
\end{itemize}

\paragraph{Security: a particularly compelling analogy.}
In the security domain, this correspondence is especially profound.
CaMeL applies operating system capability-security principles to LLM agents, creating a protection layer independent of the LLM itself~\cite{debenedetti2025camel}.
IronClaw argues for protecting AI agents the way an operating system protects its processes~\cite{ironclaw2025}.
Systems security researchers have begun using the term ``agentic computing'' to define the foundational problems of agent security~\cite{agenticsecurity2025}.
These efforts underscore a critical insight: \textbf{security boundaries are primarily a ``runtime structure'' problem and only secondarily a ``prompt'' problem}, just as OS security depends first on the permission model and isolation mechanisms, and only then on application code quality.

%% ---------- Tool Bus: I/O System ----------
\subsection{Tool Bus and Agent Interconnection}
\label{sec:io}

\subsubsection{The Counterpart in LLM-Based Systems}

\textit{Classical counterpart---the I/O system:} it evolved from per-peripheral dedicated interfaces (1960s) through standardized buses (PCI/PCIe, USB; 1990s) to a unified network stack (TCP/IP); the central lesson is that \textbf{a standardized bus protocol is a prerequisite for ecosystem prosperity}.

In model-native computing, the I/O layer is undergoing a similar transition from point-to-point connections to a standardized bus.

\paragraph{Early phase (method-level):}
ReAct interleaved reasoning with acting~\cite{yao2023react}; Toolformer taught models when to call APIs~\cite{schick2023toolformer}; HuggingGPT placed an LLM in a ``controller'' role to orchestrate heterogeneous models~\cite{hugginggpt2023}; Gorilla focused on the accuracy and documentation-refresh resilience of API calls~\cite{patil2024gorilla}.
These efforts correspond to the ``direct device manipulation'' phase of early computing: each tool required a bespoke interface and calling convention, with no unified protocol standard.

\paragraph{Standardization phase (protocol-level):}
MCP (Model Context Protocol) standardizes the connection between AI applications and external data sources, tools, and workflows~\cite{mcp_intro_2026}, playing a role analogous to PCIe or USB in traditional computing: it defines unified discovery, connection, invocation, and error-handling mechanisms so that new tools can be ``plug-and-play.''
Google's A2A (Agent-to-Agent) protocol targets agent-to-agent interoperation~\cite{google_a2a_2025}, supporting capability discovery, task lifecycle management, and multimodal collaboration. It is analogous to TCP/IP in the networking stack, enabling agents from different vendors to discover and cooperate with one another.

Taken together, MCP serves as a \textbf{vertical bus} (agent$\to$tool) while A2A serves as a \textbf{horizontal interconnect} (agent$\to$agent).
As one survey observes, MCP, A2A, and related protocols such as ACP are forming the infrastructure layer for agent interoperability~\cite{agentprotocols2025}.

\subsubsection{Why the Analogy Holds}

\paragraph{Taming interface diversity.}
Before PCIe unified motherboard-level device interfaces, every device needed a dedicated driver; before MCP unified tool interfaces, every tool required bespoke API adapter code.
Both reduce system integration costs by defining a standard protocol that hides implementation differences.

\paragraph{A shared layered stack.}
Traditional I/O spans the physical layer (electrical signaling), the link layer (PCIe transaction layer), the protocol layer (NVMe/SCSI), and the application layer (filesystem).
Agent I/O follows a similar stack: transport layer (HTTP/SSE), protocol layer (MCP/A2A message formats), semantic layer (tool schemas and agent capability descriptions), and application layer (concrete tool invocations or agent collaborations).

\paragraph{Discovery and negotiation.}
A USB device, upon insertion, must be ``enumerated'': the host queries the device type, capabilities, and driver requirements.
An MCP server, upon connection, similarly undergoes ``capability discovery'': the agent queries which tools the server supports, along with each tool's parameter schema and side-effect description.

\subsubsection{Essential Differences}

\paragraph{Non-deterministic, asynchronous, side-effecting invocation.}
A traditional system call such as \texttt{read(fd, buffer, size)} is fully predictable: identical parameters yield identical results.
Tool invocations, by contrast, may be non-deterministic (search engine results change over time), asynchronous (long-running tasks require callbacks), and side-effecting (sending email, executing code).

\paragraph{Far more complex error handling.}
PCIe devices either respond correctly or return a well-defined error code; MCP tools may return semantically ambiguous errors (``the result is uncertain'') or---worse---results that appear correct but are in fact wrong.
This compels the agent I/O layer to introduce \textbf{verification and redundancy} mechanisms: critical operations require result validation, analogous to quorum reads in distributed systems.

%% ---------- Multi-Agent: Distributed System ----------
\subsection{Multi-Agent Collaboration: Distributed Systems with Semantics}
\label{sec:distributed}

\subsubsection{The Counterpart in LLM-Based Systems}

\textit{Classical counterpart---distributed systems:} cooperation across networked machines confronting task decomposition, state consistency, fault recovery, deadlock avoidance, and load balancing~\cite{hennessy2017quantitative}.

When multiple agents collaborate on a complex task, the system exhibits characteristics of distributed computing.

AutoGen treats multi-agent conversation as general infrastructure~\cite{wu2023autogen}, where multiple agents cooperate through message passing, much like distributed nodes.
MetaGPT uses pipeline-style Standard Operating Procedures to make multi-agent collaboration explicitly procedural~\cite{metagpt2023}, analogous to industrial-pipeline task decomposition, where each agent handles one stage (requirements analysis $\to$ design $\to$ coding $\to$ testing).
Internet of Agents imagines the agent ecosystem as an ``internet'' rather than a single machine~\cite{ioa2024}, emphasizing agent integration protocols, dynamic team formation, and distributed environments.

From an architectural perspective, clear correspondences emerge:
a single agent maps to a single process, multi-agent collaboration to multithreading or multiprocessing, and a cross-node agent cluster to a distributed system.
This immediately introduces the classical distributed-system challenges: task decomposition and scheduling, state synchronization and consistency, fault detection and recovery, deadlock avoidance, and load balancing.

\subsubsection{Why the Analogy Holds}

\paragraph{Communication overhead versus parallelism benefit.}
In distributed systems, Amdahl's law tells us that parallel speedup is limited by the serial fraction.
In multi-agent systems, Heuristic~III (the agent speedup heuristic) describes analogous behavior: when orchestration efficiency $E$ is insufficient, simply adding more agents yields diminishing returns.
AutoGen's multi-agent experiments are consistent with this picture: increasing the agent count from 2 to 8 produced progressively smaller reductions in task completion time~\cite{wu2023autogen}.
Industrial case studies qualitatively illustrate that coordination topology is itself a design variable: Fountain adopted a hub-and-spoke topology (a central Copilot coordinating sub-agents for screening, document generation, and sentiment analysis)~\cite{anthropic2026agenticreport}; MetaGPT employs a pipeline topology (SOP-driven sequential collaboration); AutoGen uses a peer-to-peer topology (conversational coordination). We treat these as qualitative illustrations of topology diversity rather than quantitative calibrations of $E$---the Fountain figures in particular are vendor-reported and non-identifiable, as discussed in Section~\ref{sec:law3}.
Different topologies yield different $E$ values across task scenarios, implying that the $E$ in Heuristic~III is not a fixed scalar but a design variable coupled to the coordination structure.

\paragraph{The need for explicit coordination mechanisms.}
Distributed systems use consensus protocols such as Raft~\cite{raft2014} to help multiple nodes agree on decisions; multi-agent systems need ``debate,'' ``voting,'' or ``hierarchical arbitration'' mechanisms for multiple agents to converge on a conclusion.

\paragraph{A consistency--availability tension.}
The CAP theorem states that a distributed system cannot simultaneously guarantee consistency, availability, and partition tolerance~\cite{cap2002}.
In multi-agent systems, an analogous tension exists between ``semantic consistency'' (whether agents share the same understanding of a fact) and ``response speed'' (whether to wait for all agents to finish reasoning before producing a conclusion).

\subsubsection{Essential Differences}

\paragraph{Lossy, low-precision communication.}
Distributed nodes communicate through structured protocols (e.g., gRPC with Protocol Buffers), yielding precise, lossless message semantics.
Agents communicate through natural language or semi-structured messages, where information may be misinterpreted, omitted, or distorted during transmission; \textit{semantic noise} is a challenge unique to agent-based distributed systems.

\paragraph{Gradual rather than binary failure.}
Distributed-node failures are typically binary (online or offline); agent ``failures'' can be \textit{gradual}: the agent has not crashed, but its reasoning quality degrades, it drifts from the task objective, or it begins to hallucinate.
Detecting ``whether an agent is executing the task correctly'' is far harder than detecting ``whether a node is alive.''

\paragraph{Emergence, not deterministic composition.}
The behavior of a distributed system is the deterministic composition of its nodes' behaviors; multi-agent systems may produce \textbf{emergent effects}, where each agent executes correctly in isolation, yet the collective outcome is unpredictable.
This places new demands on system debuggability and accountability.

%% ---------- Architecture Diagram and Control Flow ----------

Figure~\ref{fig:arch} presents the proposed layered architecture for model-native computing.
Each layer corresponds to one of the component levels analyzed in this section: the model core (CPU) $\to$ KV cache (caches) $\to$ memory tier (virtual memory) $\to$ agent/OS layer (operating system) $\to$ I/O and persistence tier (peripherals) $\to$ distributed substrate (cluster).
Note the bypass arrows from the agent/OS layer directly to the memory tier and I/O tier, reflecting the dual-plane architecture's assignment of resource-management responsibility to the deterministic control plane.

\begin{figure}[htbp]
\centering
\resizebox{0.97\textwidth}{!}{%
\begin{tikzpicture}[
    node distance=4mm and 5mm,
    layer/.style={
        draw, rounded corners=4pt, minimum width=9.6cm,
        minimum height=1cm, font=\small, align=center,
        line width=0.7pt
    },
    subbox/.style={
        draw, rounded corners=3pt, minimum height=0.85cm,
        font=\footnotesize, align=center, line width=0.6pt
    },
    arrow/.style={-{Latex[length=2mm]}, thick, black!40},
    bypass/.style={
        -{Latex[length=2.8mm]}, line width=2pt,
        modelnative!85!black
    },
]

% ============================================================
% Background: Probabilistic Execution Plane  (L1--L3)
% ============================================================
\fill[classical!10, rounded corners=8pt]
    (-6.4,-6.5) rectangle (6.4,0.0);

% Background: Deterministic Control Plane  (L4--L5)
\fill[modelnative!10, rounded corners=8pt]
    (-6.4,0.6) rectangle (6.4,5.3);

% ============================================================
% L6 -- Application Layer (above both planes)
% ============================================================
\node[layer, fill=gray!18, draw=gray!55]
    (L6) at (0,6.2)
    {\textbf{L6 \, Application Layer}\\[-1pt]
     {\scriptsize User tasks \, agent apps \, workflow automation \, governance}};

% ============================================================
% L5 -- Orchestration Layer
% ============================================================
\node[layer, fill=modelnative!20, draw=modelnative!65]
    (L5) at (0,4.4)
    {\textbf{L5 \, Orchestration Layer}\\[-1pt]
     {\scriptsize Planning $\cdot$ Agent scheduling $\cdot$ Permission / Approval $\cdot$
      Failure recovery $\cdot$ State commitment}};

% ============================================================
% L4 -- Semantic Interface Layer
% ============================================================
\node[layer, fill=modelnative!20, draw=modelnative!65]
    (L4) at (0,2.8)
    {\textbf{L4 \, Semantic Interface Layer}\\[-1pt]
     {\scriptsize Tool schemas (Tool ABI) $\cdot$ Memory APIs $\cdot$ Trace APIs $\cdot$
      MCP / A2A protocols}};

% ---- Dashed interface line at L3/L4 boundary ----
\draw[dashed, thick, black!35, line width=0.9pt]
    (-6.4,0.3) -- (6.4,0.3);
\node[font=\tiny\itshape, text=black!45, anchor=west] at (-6.3,0.3)
    {L3--L4 interface contract};

% ============================================================
% L3 -- Context Management Layer
% ============================================================
\node[layer, fill=classical!20, draw=classical!60]
    (L3) at (0,-0.9)
    {\textbf{L3 \, Context Management Layer}\\[-1pt]
     {\scriptsize Hot / warm / cold tiering $\cdot$ Context compilation $\cdot$
      Summarization \& retrieval $\cdot$ Virtual context}};

% ============================================================
% L2 -- Inference Serving Layer (split into sub-components)
% ============================================================
\node[subbox, fill=classical!20, draw=classical!60,
      minimum width=4.4cm]
    (L2core) at (-2.2,-2.7)
    {\textbf{Model Core}\\[-1pt]
     {\scriptsize Foundation model / Inference strategy}};

\node[subbox, fill=classical!20, draw=classical!60,
      minimum width=2.3cm]
    (L2sess) at (1.5,-2.35)
    {\textbf{Session KV}\\[-1pt]
     {\scriptsize Per-request cache}};

\node[subbox, fill=classical!20, draw=classical!60,
      minimum width=2.3cm]
    (L2pref) at (1.5,-3.05)
    {\textbf{Prefix / Shared KV}\\[-1pt]
     {\scriptsize Cross-request reuse}};

% thin bracket grouping L2 sub-boxes
\draw[classical!50, line width=0.5pt]
    (3.0,-2.1) -- (3.25,-2.1) -- (3.25,-3.3) -- (3.0,-3.3);
\node[font=\tiny\bfseries, text=classical!70!black, anchor=west]
    at (3.35,-2.7) {L2};

% ============================================================
% L1 -- Physical Execution Layer
% ============================================================
\node[layer, fill=classical!20, draw=classical!60]
    (L1) at (0,-4.7)
    {\textbf{L1 \, Physical Execution Layer}\\[-1pt]
     {\scriptsize GPU / TPU / ASIC \, matrix ops $\cdot$ Attention kernels $\cdot$
      Sampling \& decoding}};

% ============================================================
% L0 -- Distributed Substrate (below both planes)
% ============================================================
\node[layer, fill=gray!18, draw=gray!55]
    (L0) at (0,-6.0)
    {\textbf{Distributed Substrate}\\[-1pt]
     {\scriptsize Continuous batching $\cdot$ Prefill / decode decoupling $\cdot$
      Parallel routing $\cdot$ Replication $\cdot$ Consistency}};

% ============================================================
% Plane labels (inside backgrounds, top-left)
% ============================================================
\node[font=\small\bfseries, text=classical!80!black, anchor=north west,
      fill=classical!6, fill opacity=0.7, text opacity=1,
      rounded corners=2pt, inner sep=2.5pt]
    at (-6.2,-0.1)
    {Probabilistic Execution Plane};

\node[font=\small\bfseries, text=modelnative!80!black, anchor=north west,
      fill=modelnative!6, fill opacity=0.7, text opacity=1,
      rounded corners=2pt, inner sep=2.5pt]
    at (-6.2,5.1)
    {Deterministic Control Plane};

% ============================================================
% Vertical arrows (inter-layer data / service flow)
% ============================================================
\draw[arrow] (L6)   -- (L5);
\draw[arrow] (L5)   -- (L4);
\draw[arrow] (L4)   -- (L3);
\draw[arrow] (L3)   -- (L2core);
\draw[arrow] (L2core) -- (L1);
\draw[arrow] (L1)   -- (L0);
\draw[arrow] (L2sess)  -- (L2core);
\draw[arrow] (L2pref)  -- (L2core);

% ============================================================
% BYPASS ARROWS -- the central visual feature
% ============================================================
% LEFT bypass: L5 down to L2 (KV cache management)
\draw[bypass]
    ([xshift=-1mm]L5.west) -- ++(-1.4,0)
    coordinate (bl1) -- ++(0,-5.9)
    coordinate (bl2) -- ++(1.4,0)
    ([xshift=-1mm]L2core.west);
\node[font=\scriptsize\bfseries, text=modelnative!85!black,
      rotate=90, anchor=south]
    at ([xshift=-8mm]bl2) {Direct KV cache control};

% RIGHT bypass: L4 down to L3 (context management) and L2
\draw[bypass]
    ([xshift=1mm]L4.east) -- ++(1.4,0)
    coordinate (br1) -- ++(0,-4.2)
    coordinate (br2) -- ++(-1.4,0)
    ([xshift=1mm]L3.east);

\draw[bypass]
    ([xshift=1mm]L5.east) -- ++(2.2,0)
    coordinate (br3) -- ++(0,-7.7)
    coordinate (br4) -- ++(-1.8,0)
    ([xshift=1mm]L1.east);
\node[font=\scriptsize\bfseries, text=modelnative!85!black,
      rotate=-90, anchor=south]
    at ([xshift=8mm]br4) {Direct inference control};

\end{tikzpicture}}
\caption{ICA layered architecture for model-native computing.
  \textbf{Top to bottom:} L6 Application, L5 Orchestration, L4 Semantic Interface
  (Deterministic Control Plane, orange), dashed L3--L4 interface boundary,
  L3 Context Management, L2 Inference Serving (Model Core + Session KV +
  Prefix/Shared KV), L1 Physical Execution (Probabilistic Execution Plane, blue),
  and Distributed Substrate.
  Thick orange bypass arrows show the Deterministic Control Plane reaching
  through the interface boundary to directly manage KV caches (L2) and
  hardware execution (L1).}
\label{fig:arch}
\end{figure}


Figure~\ref{fig:flow} illustrates the control flow of a typical request.
The flow embodies the interaction pattern of the dual-plane architecture: the probabilistic plane is responsible for generation and understanding, while the deterministic plane handles scheduling, approval, and state management.

\begin{figure}[htbp]
\centering
\resizebox{0.95\textwidth}{!}{%
\begin{tikzpicture}[
    % ---- Styles ----
    procB/.style={draw=classical!60, rounded corners=4pt,
        minimum width=6.2cm, minimum height=9mm, align=center,
        fill=classical!15, font=\small, line width=0.65pt},
    procO/.style={draw=modelnative!60, rounded corners=4pt,
        minimum width=6.2cm, minimum height=9mm, align=center,
        fill=modelnative!15, font=\small, line width=0.65pt},
    procN/.style={draw=gray!50, rounded corners=4pt,
        minimum height=9mm, align=center,
        fill=gray!10, font=\small, line width=0.65pt},
    dec/.style={draw, diamond, aspect=2.6, align=center,
        inner sep=2pt, font=\small, line width=0.65pt},
    arrB/.style={-{Latex[length=2.2mm]}, thick, classical!75},
    arrO/.style={-{Latex[length=2.2mm]}, thick, modelnative!75},
    arrG/.style={-{Latex[length=2.2mm]}, thick, gray!55},
    hdr/.style={font=\footnotesize\bfseries, inner sep=3pt,
        rounded corners=2pt, fill opacity=0.55, text opacity=1},
    node distance=5mm and 5mm,
]

% ============================================================
%  GEOMETRY
% ============================================================
\pgfmathsetmacro{\cxL}{-3.8}   % left-lane centre x
\pgfmathsetmacro{\cxR}{+3.8}   % right-lane centre x
\pgfmathsetmacro{\halfW}{3.5}  % half-width of each lane band

\def\yBot{-10.2}
\def\yTop{1.35}

% ============================================================
%  SWIM-LANE BACKGROUNDS  (blue left, orange right)
% ============================================================
\fill[classical!6, rounded corners=8pt]
    (\cxL-\halfW, \yBot-0.35) rectangle
    (\cxL+\halfW, \yTop+0.75);

\fill[modelnative!6, rounded corners=8pt]
    (\cxR-\halfW, \yBot-0.35) rectangle
    (\cxR+\halfW, \yTop+0.75);

% ---- Lane header pills ----
\node[hdr, text=classical!80!black, fill=classical!10]
    at (\cxL, \yTop+0.45) {Probabilistic Execution};
\node[hdr, text=modelnative!80!black, fill=modelnative!10]
    at (\cxR, \yTop+0.45) {Deterministic Control};

% ============================================================
%  INPUT  (spans both lanes, centred)
% ============================================================
\node[procN, minimum width=13.0cm]
    (req) at (0, 0.75)
    {User request arrives: goal, constraints, environment handle};

% ============================================================
%  LEFT LANE -- Probabilistic Execution  (blue)
% ============================================================
\node[procB] (plan) at (\cxL, 0.0)
    {Agent kernel parses task and loads\\base policies, project memory};

\node[procB, below=4.5mm of plan]
    (perm) {Permission \& capability check\\(deterministic gate)};

\node[procB, below=4.5mm of perm]
    (prefill) {Context compilation \& address resolution:\\
    hot-window loading, summary backfill,\\
    prefix-cache hit, prefill/decode scheduling};

\node[procB, below=4.5mm of prefill]
    (gen) {Model execution \& control loop:\\
    generate action, validate result,\\
    retry or fork if needed};

% ============================================================
%  RIGHT LANE -- Deterministic Control  (orange)
% ============================================================
\node[procO]
    (policy) at (\cxR, 0.0)
    {Load permission config,\\
    tool-capability annotations};

\node[dec, below=6mm of policy,
      draw=classical!50, fill=modelnative!8]
    (tool) {Tool / sub-agent\\needed?};

% Split-colour diagonal lines inside diamond
\draw[modelnative!40, line width=0.4pt] (tool.south west) -- (tool.north east);
\draw[classical!40, line width=0.4pt]   (tool.north west) -- (tool.south east);

\node[procO, below=5mm of tool]
    (syscall) {System call path:\\
    MCP / tool invocation, sandboxed\\
    execution, result return};

\node[procO, below=4.5mm of syscall]
    (commit) {State commit: update short-term memory,\\
    long-term memory, execution log \&\\
    auditable evidence};

% ============================================================
%  CONVERGENCE / OUTPUT  (spans both lanes)
% ============================================================
\node[procN, minimum width=13.0cm]
    (resp) at (0, \yBot)
    {Output response or produce committable artifact};

% ---- Convergence dot ----
\fill[gray!30] (0, \yBot+0.55) circle (3pt);
\node[font=\tiny, text=gray!50, right=3pt] at (0.05, \yBot+0.55) {merge};

% ============================================================
%  ARROWS
% ============================================================
% Input splits into both lanes
\coordinate (split) at (0, 0.15);
\draw[arrG] (req.south) -- (split);
\draw[arrG] (split) -| (plan.north);
\draw[arrG] (split) -| (policy.north);

% Left lane (probabilistic)
\draw[arrB] (plan) -- (perm);
\draw[arrB] (perm) -- (prefill);
\draw[arrB] (prefill) -- (gen);

% Right lane (deterministic)
\draw[arrO] (policy) -- (tool);
\draw[arrO] (tool.south) --
    node[right, font=\scriptsize, text=modelnative!70!black] {Yes}
    (syscall);
\draw[arrO] (syscall) -- (commit);

% Cross-lane: decision "No" feeds context compilation
\draw[arrB] (tool.west) --
    node[above, font=\scriptsize, text=classical!70!black] {No}
    (prefill.east);

% Cross-lane: model generation feeds state commit
\draw[arrO] (gen.east) -- (commit.west);

% Convergence to output
\coordinate (merge) at (0, \yBot+0.9);
\draw[arrG] (gen.south)  |- (merge);
\draw[arrG] (commit.south) |- (merge);
\draw[arrG] (merge) -- (resp.north);

\end{tikzpicture}}
\caption{Request flow through the model-native computing architecture as a dual-plane swim-lane diagram.
Blue nodes belong to the probabilistic execution plane; orange nodes belong to the deterministic control plane.
The split-coloured diamond represents the control-plane decision that gates both paths.
Grey boxes span both planes.}
\label{fig:flow}
\end{figure}

